% 1. Поиск в PDF и кодировки (cmap строго перед fontenc!)
\usepackage{cmap}
\usepackage[T2A]{fontenc}
\usepackage[utf8]{inputenc}
\usepackage[russian]{babel}
\usepackage{cm-super} % векторные Type 1 шрифты для печати без пропусков

% 2. Поля по ГОСТ 7.32
\usepackage[left=30mm, right=15mm, top=20mm, bottom=20mm]{geometry}

% 3. Межстрочный интервал и абзацы
\usepackage{setspace}
\onehalfspacing
\usepackage{indentfirst}
\setlength{\parindent}{1.25cm}

% 4. Математика
\usepackage{amsmath, amsfonts, amssymb}

% 5. Таблицы и графика
\usepackage{graphicx}
\graphicspath{{exports_fig/}}
\usepackage{booktabs}
\usepackage{array}
\usepackage{float}
\usepackage{multirow}

% 6. Подписи к рисункам и таблицам
\usepackage{caption}
\DeclareCaptionLabelSeparator{endash}{\space---\space}
\captionsetup[figure]{labelsep=endash, justification=centering, font=small}
\captionsetup[table]{labelsep=endash, justification=raggedright, singlelinecheck=false, font=small}

% 7. Заголовки разделов (titlesec)
\usepackage{titlesec}
\titleformat{\section}[hang]
    {\raggedright\bfseries\large}
    {\hspace*{\parindent}\thesection}
    {0.5em}
    {}

\titleformat{\subsection}[hang]
    {\raggedright\bfseries\normalsize}
    {\hspace*{\parindent}\thesubsection}
    {0.5em}
    {}

\titlespacing*{\section}{0pt}{14pt plus 2pt minus 2pt}{8pt plus 2pt minus 2pt}
\titlespacing*{\subsection}{0pt}{10pt plus 2pt minus 2pt}{6pt plus 2pt minus 2pt}

% 8. Оглавление (опция [titles] устраняет конфликт с titlesec)
\usepackage[titles]{tocloft}
\renewcommand{\contentsname}{СОДЕРЖАНИЕ}
\renewcommand{\cfttoctitlefont}{\hfill\bfseries\normalsize}
\renewcommand{\cftaftertoctitle}{\hfill\null\smallskip}
\renewcommand{\cftsecleader}{\cftdotfill{\cftdotsep}}
\renewcommand{\cftsecfont}{\normalfont}
\renewcommand{\cftsecpagefont}{\normalfont}
\setlength{\cftbeforesecskip}{0.3em}

% 9. Подсветка кода (listings)
\usepackage{xcolor}
\usepackage{listings}
\definecolor{codegreen}{rgb}{0,0.6,0}
\definecolor{codegray}{rgb}{0.5,0.5,0.5}
\definecolor{codepurple}{rgb}{0.58,0,0.82}
\definecolor{backcolour}{rgb}{0.96,0.96,0.96}

\lstdefinestyle{gostcode}{
    backgroundcolor=\color{backcolour},
    commentstyle=\color{codegreen},
    keywordstyle=\color{magenta},
    numberstyle=\tiny\color{codegray},
    stringstyle=\color{codepurple},
    basicstyle=\ttfamily\footnotesize,
    breaklines=true,
    captionpos=t,
    numbers=left,
    numbersep=5pt,
    tabsize=4,
    frame=single,
    keepspaces=true,
    extendedchars=true
}
\lstset{style=gostcode}

% 10. Ссылки (грузятся почти в самом конце)
\usepackage{hyperref}
\hypersetup{
    colorlinks=true,
    linkcolor=black,
    filecolor=black,
    urlcolor=blue,
    citecolor=black
}
